\chapter{プロンプトテンプレート集}
\label{app:prompts}

本付録では、本書で紹介したプロンプトのエッセンスをテンプレートとして整理する。各テンプレートには、ROCF（Role, Output, Context, Format）の各要素がどこに対応するかを注釈として示している。

\begin{itembox}[l]{ROCF要素の凡例}
\begin{itemize}
\item \textbf{【R】}Role（役割）: AIに与える役割設定
\item \textbf{【O】}Output（出力）: 求める出力の内容
\item \textbf{【C】}Context（文脈）: 背景情報や条件
\item \textbf{【F】}Format（形式）: 出力の形式指定
\end{itemize}
\end{itembox}

\section{文献調査用プロンプト}

\begin{promptbox}
\textbf{テンプレート1: サーベイの起点を見つける}

【R】あなたは[分野名]の研究者で、文献調査の専門家です。

【O】以下のテーマに関するサーベイ論文・レビュー論文を\\
探すための検索戦略を提案してください。

【C】テーマ: [テーマ] / 知りたいこと: [具体的な問い] /\\
すでに知っていること: [既知の情報]

【F】推奨する検索キーワード（英語）5セット /\\
推奨する学術データベース / 注意すべき関連分野
\end{promptbox}

\begin{promptbox}
\textbf{テンプレート2: 論文の要約を依頼する}

【R】あなたは学術論文の読解に精通した研究アシスタントです。

【O】以下の論文情報を、構造化された要約にまとめて\\
ください。

【C】タイトル: [タイトル] / 著者: [著者] /\\
概要: [Abstractまたは要点]

【F】研究の目的（2文以内）/ 手法の概要（3文以内）/\\
主な結果（箇条書き3点）/ 限界・今後の課題（2点）/\\
自分の研究との関連（※ここは自分で記入すること）
\end{promptbox}

\begin{promptbox}
\textbf{テンプレート3: 研究動向の整理}

【R】あなたは[分野名]の動向に詳しい研究者です。

【O】以下のキーワードに関する研究動向を整理してください。

【C】キーワード: [キーワード] / 期間: 過去[5]年間 /\\
対象: [学会名/ジャーナル名]

【F】主要な研究潮流（3〜5つ）/ 最近のトレンド・転換点 /\\
未解決の課題・今後の方向性\\
※ 挙げられた論文は必ず実在を確認すること
\end{promptbox}

\begin{promptbox}
\textbf{テンプレート4: 論文間の関係性の整理}

【R】あなたは文献レビューの専門家です。

【O】以下の複数の論文の関係性を整理してください。

【C】論文1: [タイトル, 著者, 年] --- [要点]\\
論文2: [タイトル, 著者, 年] --- [要点]\\
論文3: [タイトル, 著者, 年] --- [要点]

【F】各論文の共通点と相違点を表形式で /\\
研究の発展の流れ（時系列）/ 未解決の問題
\end{promptbox}

\begin{promptbox}
\textbf{テンプレート5: 批判的読解の支援}

【R】あなたは厳格な査読者です。

【O】以下の論文の主張について、批判的な視点からの質問を\\
生成してください。

【C】論文の主張: [主張] / 使用された手法: [手法] /\\
実験の規模: [データ数等]

【F】以下の観点から各3問ずつ: 方法論的な妥当性 /\\
結果の解釈 / 一般化可能性
\end{promptbox}

\section{英語論文用プロンプト}

\begin{promptbox}
\textbf{テンプレート1: Abstractの英訳・改善}

【R】あなたは学術英語ライティングの専門家です。\\{}[分野名]の論文に精通しています。

【O】以下の日本語要旨を、学術英語のAbstractとして\\
書き直してください。

【C】[日本語の要旨] / 対象ジャーナル/学会: [名称] /\\
字数制限: [語数]語以内

【F】構成: Background → Purpose → Methods →\\
Results → Conclusion / 学術英語の慣習に従った表現 /\\
出力後、使用した主要な学術表現のリストを付記
\end{promptbox}

\begin{promptbox}
\textbf{テンプレート2: Introduction の段落構成}

【R】あなたは学術論文のライティングコーチです。

【O】以下の情報をもとに、Introductionの段落構成案を\\
作成してください。

【C】研究分野: [分野] / 研究の背景: [2-3文] /\\
既存研究の課題: [箇条書き] / 本研究の提案: [1-2文]

【F】各段落の役割 / 使うべき英語表現の候補（2-3個）/\\
段落間の接続の仕方
\end{promptbox}

\begin{promptbox}
\textbf{テンプレート3: 英文の校正・改善}

【R】あなたはネイティブの学術英語校正者です。

【O】以下の英文を校正し、より学術的に洗練された表現に\\
改善してください。

【C】[校正対象の英文] / 論文のセクション:\\
{[}Introduction / Methods / Results / Discussion] /\\
分野: [分野名]

【F】修正箇所を明示（元の表現→修正後の表現）/\\
各修正の理由を簡潔に説明 / 全体的な改善提案
\end{promptbox}

\begin{promptbox}
\textbf{テンプレート4: 査読コメントへの回答文作成}

【R】あなたは学術論文の査読対応に精通した研究者です。

【O】以下の査読コメントに対する回答文の下書きを\\
作成してください。

【C】査読コメント: [コメント] /\\
対応方針: [どのように対応するか]

【F】英語で記述 / ``We thank the reviewer for...'' で\\
始める丁寧な書き出し / 論文の修正箇所を明示
\end{promptbox}

\begin{promptbox}
\textbf{テンプレート5: 関連研究セクションの構成}

【R】あなたは学術論文の Related Work セクションの\\
構成アドバイザーです。

【O】以下の先行研究リストを、Related Workセクション\\
として効果的に構成する方法を提案してください。

【C】[先行研究のリスト] / 本研究のテーマ: [テーマ]

【F】サブセクションの分け方 /\\
各サブセクションで述べるべき内容 /\\
先行研究から本研究への橋渡しの書き方
\end{promptbox}

\section{コード生成用プロンプト}

\begin{promptbox}
\textbf{テンプレート1: 分析コードの雛形生成}

【R】あなたはPythonのデータ分析に精通したプログラマーです。

【O】以下の分析タスクを実行するPythonコードを\\
生成してください。

【C】目的: [例: CSVファイルからデータを読み込み、\\
基本統計量を算出] / データの形式: [例: CSV、列名はA, B, C] /\\
使用ライブラリ: [例: pandas, matplotlib]

【F】コメント付きの読みやすいコード / 各ステップの説明 /\\
エラーハンドリングを含む
\end{promptbox}

\begin{promptbox}
\textbf{テンプレート2: 既存コードのリファクタリング}

【R】あなたはコードレビューの専門家です。

【O】以下のコードをリファクタリングしてください。

【C】[リファクタリング対象のコード] /\\
問題点: [例: 可読性が低い、重複が多い]

【F】リファクタリング後のコード /\\
変更点のリスト（何を、なぜ変更したか）
\end{promptbox}

\begin{promptbox}
\textbf{テンプレート3: テストコードの生成}

【R】あなたはソフトウェアテストの専門家です。

【O】以下の関数に対するユニットテストを生成してください。

【C】[テスト対象の関数] / テストフレームワーク: [例: pytest]

【F】正常系テスト（3ケース以上）/\\
異常系テスト（2ケース以上）/\\
エッジケーステスト（2ケース以上）
\end{promptbox}

\begin{promptbox}
\textbf{テンプレート4: エラーのデバッグ支援}

【R】あなたはPythonのデバッグに精通したエンジニアです。

【O】以下のエラーの原因と修正方法を教えてください。

【C】コード: [エラーが発生するコード] /\\
エラーメッセージ: [エラーメッセージ] /\\
実行環境: [例: Python 3.11, macOS]

【F】エラーの原因の説明（初学者にも分かるように）/\\
修正後のコード / 同様のエラーを防ぐためのヒント
\end{promptbox}

\begin{promptbox}
\textbf{テンプレート5: README.mdの生成}

【R】あなたはオープンソースプロジェクトのドキュメント専門家です。

【O】以下のプロジェクト情報から、GitHub用の\\
README.mdを生成してください。

【C】プロジェクト名: [名前] / 概要: [1-2文] /\\
使用言語: [例: Python] / 主な機能: [箇条書き]

【F】英語で記述 / 構成: Title → Badge → Overview →\\
Installation → Usage → Project Structure → License
\end{promptbox}

\section{データ分析用プロンプト}

\begin{promptbox}
\textbf{テンプレート1: 分析計画の立案}

【R】あなたはデータサイエンティストです。

【O】以下のデータと研究目的に対する分析計画を\\
提案してください。

【C】研究目的: [目的] /\\
データの概要: [データの種類、サイズ、変数] /\\
仮説: [検証したい仮説]

【F】推奨する分析手法（3つ以上）と理由 /\\
分析の手順（ステップバイステップ）/ 必要な前処理 /\\
結果の解釈の仕方
\end{promptbox}

\begin{promptbox}
\textbf{テンプレート2: グラフの改善提案}

【R】あなたはデータ可視化の専門家です。

【O】以下のグラフの改善提案をしてください。

【C】グラフの種類: [例: 棒グラフ] /\\
何を示したいか: [意図] / 現在の問題: [例: 見づらい] /\\
使用ツール: [例: matplotlib]

【F】改善ポイント（5つ以上）/ 改善後のコード例 /\\
なぜその改善が効果的かの説明
\end{promptbox}

\begin{promptbox}
\textbf{テンプレート3: 統計的検定の選択}

【R】あなたは統計学の専門家です。

【O】以下の研究状況に適した統計的検定手法を\\
推薦してください。

【C】比較したいもの: [例: 2群の平均値の差] /\\
データの性質: [例: 正規分布かどうか、対応の有無] /\\
有意水準: [例: 0.05]

【F】推薦する検定手法とその理由 / 前提条件の確認方法 /\\
Pythonでの実装コード / 結果の解釈の仕方 /\\
論文での記載例（英語）
\end{promptbox}

\section{スライド作成用プロンプト}

\begin{promptbox}
\textbf{テンプレート1: スライド構成の立案}

【R】あなたは学会発表のプレゼンテーションコーチです。

【O】以下の研究内容を、[発表時間]分の学会発表\\
スライドの構成案にしてください。

【C】研究タイトル: [タイトル] / 対象学会: [学会名] /\\
聴衆: [同分野の研究者/広い分野/実務者] /\\
研究の概要: [3-5文]

【F】各スライドについて: スライドタイトル /\\
含めるべき内容 / 推定所要時間 / 注意点
\end{promptbox}

\begin{promptbox}
\textbf{テンプレート2: スライドの文字量削減}

【R】あなたはプレゼンテーションデザインの専門家です。

【O】以下のスライドの内容を、視覚的に見やすくなるよう\\
文字量を削減してください。

【C】現在のスライド内容: [テキスト] / 最も伝えたいこと: [1文]

【F】削減後のテキスト（元の50\%以下の文字数）/\\
キーワードの強調提案 / 図表化できる部分の提案
\end{promptbox}

\begin{promptbox}
\textbf{テンプレート3: Q\&A想定質問の準備}

【R】あなたは[分野名]の学会に参加する研究者です。

【O】以下の発表内容に対して、Q\&Aセッションで出そうな\\
質問を生成してください。

【C】発表タイトル: [タイトル] / 発表の概要: [3-5文] /\\
使用した手法: [手法] / 主な結果: [結果]

【F】以下のカテゴリに分けて各3問:\\
手法に関する質問 / 結果の解釈に関する質問 /\\
今後の展望に関する質問 / 想定される厳しい質問\\
各質問に回答のポイントを付記
\end{promptbox}

\section{ガクチカ・ES用プロンプト}

\begin{promptbox}
\textbf{テンプレート1: 経験の棚卸し}

【R】あなたは大学院生専門のキャリアカウンセラーです。

【O】以下の情報をもとに、ガクチカの素材を引き出す\\
ための深掘り質問を10個生成してください。

【C】専攻: [専攻] / 研究テーマ: [テーマ] /\\
研究以外の活動: [箇条書き] /\\
大学院生活で苦労したこと: [箇条書き]

【F】各質問に対して: 質問文 /\\
この質問が引き出そうとしているポイント / 回答のヒント
\end{promptbox}

\begin{promptbox}
\textbf{テンプレート2: STAR法での構造化}

【R】あなたはESライティングの専門コーチです。

【O】以下の経験メモを、STAR法で構造化してください。

【C】[経験の箇条書きメモ] / 志望業界: [業界]

【F】S（状況）: [3文以内] / T（課題）: [2文以内] /\\
A（行動）: [3-4文] /\\
R（結果）: [2文以内、定量的な成果を含む]\\
＋各要素の改善ポイント
\end{promptbox}

\begin{promptbox}
\textbf{テンプレート3: 採用担当者視点での添削}

【R】あなたは[業界名]の採用担当者で、\\
年間500本以上のESを読んでいます。

【O】以下のガクチカを、7つの観点で評価・添削して\\
ください。

【C】[ガクチカの文章] / 志望職種: [職種]

【F】以下の各観点について ○/△/× と改善提案:\\
(1) 第一印象 (2) 課題の明確さ (3) 行動の具体性\\
(4) 主体性の伝わり方 (5) 成果の定量性\\
(6) 入社後の再現性 (7) 全体の読みやすさ\\
＋修正案（修正箇所を明示）
\end{promptbox}
